\subsection{Phonology of Filipino}

Filipino phonology is relatively simple in its native structure;
traditionally, it is described as having a basic three-vowel system
\textit{(/i, u, a/)}~\cite{kaufman2010} that undergoes allophonic
lowering of high vowels to [E] and [O]~\cite{llamzon1966} in certain
prosodic domains, while consonants include voiceless stops \textit{(/p, t, k, \textipa{P}/)},
voiced stops \textit{(/b, d, g/)}, spirants \textit{(/s, h/)}, nasals
\textit{(/m, n, \textipa{N}/)}, a lateral \textit{(/l/)}, and a flap/tap
\textit{(/r/)}. Due to the borrowing of Spanish and English loanwords,
its modern structure has evolved in complexity to adapt for the growing 
phonemes.

\textbf{Phoneme expansion}

Due to the introduction of loanwords, phonemes structures have evolved
to accommodate the growing library of loanwords being adapted to Filipino,
such examples being phonemes like affricates \textit{/f, v, z/}, and fricatives 
\textit{/\textipa{S}, \textipa{dZ}, \textipa{tS}/} that entered the language 
primarily through loanwords. The introduction of new palatal variants create 
adapted variations of these phonemes to match local phonetics, such that 
\textit{/\textipa{tS}/ and /\textipa{dZ}/}, are adapted into the language as 
\textit{/ts/ and /dy/}~\cite{llamzon1966}, or are integrated as full phonemes 
within the existing Filipino phonetic structure~\cite{kaufman2010}.While bilingual 
speakers may pronounce the affricates or fricatives true to its given pronunciation, 
these sounds are often subject to varying degrees of nativization depending on the speaker~\cite{kaufman2010}. 

\textbf{Sociolectal variation}

Variations in Philippine English (PhE) are influenced by factors such as education, 
socioeconomic status, and exposure to English media. These variations, formally referred 
to as sociolects, are categorized in descending order of their proximity to the target 
language as acrolects, mesolects, and basilects~\cite{tayao2004}.

Both vowel inventories of PhE acrolect and mesolect feature all the vowels within the General 
American English (gAmE) vowels, 
though the mesolect shows more of a gradient, such as \textit{/\textipa{I}$\sim$i/},
\textit{/\textipa{E}$\sim$e/}, and \textit{/\textipa{A}$\sim$\textipa{\ae}/}~\cite{magpale_segmental_2024}. 
In contrast, 
the PhE basilect tends to shorten long vowels. This reduced 
inventory often leads to the merger of distinct gAmE phonemes; specifically, 
\textit{/\textipa{\ae}/} and \textit{/\textipa{V}/} are absent in the basilect inventory, but
they are both 
realized as \textit{[\textipa{A}]} (e.g., ``bad'' and ``but'' 
both surfacing as \textit{[b\textipa{A}d]} and \textit{[b\textipa{A}t]})~\cite{magpale_segmental_2024}.

\textbf{Syllabic structures}

Filipino has a basic syllable structure of consonant-vowel-consonant (CVC), but has 
evolved in complexity due to introductions of loanwords (\textit{ex.~komiks (CVCC), trump, (CCVCC)}).
While adaptation of loanwords into a Filipino allow it to mostly retain its original syllable 
structure, Filipino speakers often use vowel insertion to repair clusters that violate 
native constraints (e.g., is-kula for ``school'') or, in mesolectal varieties, deletion of the 
This vowel insert is evident in the realization of $sC$ clusters within the mesolects and basilects. 
In the mesolect, $sC$ clusters combinations like 
/sm/, /sn/, and /sl/ show occasional vowel insertion. Meanwhile the basilect undergoes the most substantial 
modification, frequently inserting a vowel \textit{/\textipa{I}/} before /s/ (e.g., 
\textit{skill} becoming \textit{[\textipa{Is.kIl}]} or \textit{sweet} 
becoming \textit{[\textipa{Is.wIt}]})~\cite{magpale_segmental_2024,magpale_hong_2024_sc}.

% idk how relevant this is, but these were also discussed in the papers
\textbf{Phonological processes}

In Filipino vocabulary, vowel length and primary stress are typically restricted to the 
final two syllables. However, loanwords place the length and stress much earlier in the word, 
such as in the Spanish loan telépono (``telephone'') or the Sanskrit loan dálità (``suffering''). 
The adaptation of loanwords also reflects the influence of the syllable-timed rhythm~\cite{dayag2007} 
of native Philippine languages on borrowed English. Mesolectal speakers often substitute the interdental 
fricatives \textit{/\textipa{T}/ and /\textipa{D}/} found in English loanwords (like think or those) with the 
alveolar stops [t] and [d], respectively. Despite these phonological shifts, these adapted forms remain highly 
intelligible to listeners familiar with the Philippine variety of English.